\headright{Berufserfahrung}

\jobtitle{DevOps-Ingenieur}{Alten (Frankreich)}{Mai 2025 – März 2026}{}
\small

\noindent \textbf{Amadeus}

\begin{itemize}[label=\tiny$\bullet$, leftmargin=1em]
	\item Migration von älteren Groovy-Pipeline-Skripten in ein internes Framework, das den gesamten Anwendungslebenszyklus abdeckt.
	\item Neuprogrammierung des Feed-Frameworks von \textbf{Bash} nach \textbf{Python} für geplante Datenlieferungen.
	\item Einführung von Branchen-Frameworks zur Automatisierung lokaler Abhängigkeiten (Localstack, Nix devenv).
	\item Entwicklung von Skripten zum dynamischen Wechseln von Maschinentypen auf EMR-Clustern von Spot- zu On-Demand-Instanzen basierend auf der Warteschlangenlatenz.
	\item Überarbeitung der internen Dokumentation zur Verbesserung der Übersichtlichkeit und Wartbarkeit.
	\item Entwurf eines HLD für das Release-Management zur Standardisierung von Deployment-Workflows gemäß den Anforderungen des Änderungsmanagements.
\end{itemize}

\vspace{3.0ex}
\jobtitle{DevOps Engineer}{DXC Technology (Bulgarien)}{02.2022 -- 04.2025} {}
\small

\noindent \textbf{Continental AG (Reifen)}
\begin{itemize}[label=\tiny$\bullet$, leftmargin=1em, nosep]
	\item Leitung eines 5-köpfigen Teams und direkte Berichterstattung über den Projektstatus an den Kunden.
	\item Optimierung der Datenverarbeitungslogik durch Migration von älteren Spark-Batch-Jobs zu Apache Flink-Streaming-Pipelines, wodurch die Latenz erheblich reduziert wurde.
	\item Beschleunigte den Projektstart durch die Behebung kritischer VPN-Verbindungsprobleme und die Automatisierung lokaler Entwicklungsprozesse.
\end{itemize}

\vspace{1em}
\noindent \textbf{Nestlé}
\begin{itemize}[label=\tiny$\bullet$, leftmargin=1em, nosep]
	\item Implementierte automatisierte strukturelle Validierungsprüfungen für Pull-Requests, um nicht konforme Bereitstellungen zu verhindern.
	\item Erweiterung der Snowflake-Dialekte in SQLFluff und automatisierte Integration in den Pre-Release-Schritt
	\item Integration von ANTLR-Parsern in die \textbf{CI/CD}-Pipeline als obligatorisches Validierungsgate vor der Veröffentlichung.
\end{itemize}

\vspace{1em}
\noindent \textbf{Ahold Delhaize}
\begin{itemize}[label=\tiny$\bullet$, leftmargin=1em]
	\item Das Testautomatisierungs-Framework wurde zu einer modularen Architektur umgestaltet, wodurch die Wiederverwendbarkeit des Codes verbessert und die Einarbeitungszeit für QA-Ingenieure verkürzt wurde.
	\item In wöchentlichen Abstimmungen mit dem QA-Team wurden operative Engpässe identifiziert, was zu einer verbesserten Teamautonomie führte.
	\item Entwurf eines HLD für den umgebungsübergreifenden Datenaustausch, das die Migration des neuen Cloud-Ingestion-Frameworks (CF2) unterstützte.
	\item Automatisierung eines Workflows zur Bereitstellung von Allure-Berichten als statische Website auf Azure Blob Storage mit benutzerdefinierter Authentifizierungs-Middleware.
	\item Leitung einer technischen Evaluierung von \textbf{SonarQube} für die automatisierte Codeanalyse, um die Integrationsanforderungen für die Plattform zu ermitteln
\end{itemize}

\vspace{1em}
\noindent \textbf{Continental AG (ADAS, Codename CAEdge)}
\begin{itemize}[label=\tiny$\bullet$, leftmargin=1em]
	\item Automatisierung einer standardisierten \textbf{VPC}-Topologie für eine AWS-Konfiguration mit mehreren Konten und deren Integration in den Account Vending Machine-Orchestrator.
	\item Entwurf und Automatisierung einer Lösung zur Cloud-Kosten-Governance für eine AWS-Organisation mit mehreren Konten, einschließlich der Implementierung von Budgetwarnungen zur Vermeidung von Überschreitungen.
	\item Verringerung der Angriffsfläche von Cloud-Anwendungen durch Behebung von CVEs, die in AWS Inspector-Berichten identifiziert wurden.
\end{itemize}

\vspace{5.0ex}
\jobtitle{Software Engineer Intern}{Sys Consulting Ltd.}{07.2021 -- 12.2021}{}
\small
Wartung von Legacy-Bankensystemen in CAVO für Windows-XP-Umgebungen.
